% New border-rank findings. Include in the checked-results part only after root review.
\section{Arithmetic replacement of the seventh-minimum certificates}

\subsection*{BR-001: an independent finite checker for all 25 prolongation cases}
Let the six pairs of $\{0,1,2,3\}$ be in lexicographic order, and number the
rectangles $q=(I,J)$ by $6\operatorname{index}(I)+\operatorname{index}(J)$.
For a rectangle subset $S$, set $A=k^S$ and let $B=C$ have the used variables
$x_{rc}$ as basis.  Its residual inclusion tensor is
\[
 T_S^\epsilon=\sum_{q=((r,s),(c,d))\in S}a_q\otimes
 (x_{rc}\otimes x_{sd}+x_{sd}\otimes x_{rc}
 +\epsilon x_{rd}\otimes x_{sc}+\epsilon x_{sc}\otimes x_{rd}),
 \qquad \epsilon\in\{-1,1\}.
\]
Here $\epsilon=-1$ is the determinant and $\epsilon=1$ the permanent.
In characteristic zero, the following 25 cases have border rank at least 14.  The listed determinant
cases are also taken with their transposes.
\[
\begin{array}{c|c|l}
\text{inherited locator}&\epsilon&S\\\hline
\text{P-09 A}&-1&0,1,3,6,7,9,12\\
\text{P-09 B}&-1&0,1,3,6,7,12,13\\
\text{P-12 C}&-1&0,1,2,3,6,7,8\\
\text{P-12 D}&-1&0,1,2,3,4,6,12\\
\text{P-14 E}&-1&0,1,2,3,6,7,9\\
\text{P-14 F}&-1&0,1,2,3,6,7,18\\
\text{P-15 det}&-1&0,1,3,6,7,9,18\\\hline
\text{P-13 1}&1&0,1,2,3,4,6,12\\
\text{P-13 2}&1&0,1,2,3,6,7,8\\
\text{P-13 3}&1&0,1,2,6,7,8,18\\
\text{P-13 4}&1&0,1,2,6,7,18,20\\
\text{P-13 5}&1&0,1,3,6,7,9,14\\
\text{P-16 1}&1&0,1,2,3,6,7,9\\
\text{P-16 2}&1&0,1,2,3,6,7,18\\
\text{P-16 3}&1&0,1,3,6,7,15,21\\
\text{P-15 per A}&1&0,1,3,6,7,9,18\\
\text{P-15 per B}&1&0,1,3,6,7,18,21\\
\text{P-19}&1&0,1,2,6,7,18,19
\end{array}
\]
The characteristic-zero values were claims of BR\_PENCIL.  The new part here
is their independent reconstruction by standard-library programs and their
arithmetic extension, used in BR-002 below.

\paragraph{Proof and complete algorithm.}
We state the characteristic-independent form of the necessary condition,
since the inherited findings stated it only in characteristic zero.
Write $U=T(C^*)\subset A\otimes B$.  If $\underline R(T)\le r$, there is
an $r$-plane $E\supset U$ such that
\[
 \dim\ker\Psi_A(E)\ge r,\qquad \dim\ker\Psi_B(E)\ge r,
\]
where $\Psi_A:E\otimes A\to B\otimes\Lambda^2A$ and
$\Psi_B:E\otimes B\to A\otimes\Lambda^2B$ are wedge maps.
For spans of $r$ independent pure $AB$ tensors, the corresponding symmetric
prolongation tensors give $r$ independent kernel vectors; contraction by a
linear form nonzero on all factors proves independence.  The source bundles
have constant rank on the Grassmannian, so these kernel-dimension conditions
are closed.  The condition $U\subset E$ is closed and the Grassmannian is
proper.  Projection therefore passes this necessary condition to all border
limits.  The open condition that the $r$ pure $AB$ tensors are independent
is dense in the $r$-summand parameter space.  This proves the assertion over
any algebraically closed field.  This is the all-field form of BR\_PENCIL
P-09's first reduction and Conner--Huang--Landsberg, Proposition 2.1.

The split torus gives $x_{rc}$ weight $R_r+C_c$ and $a_q$ weight
$-\sum_{r\in I}R_r-\sum_{c\in J}C_c$.  A separating cocharacter limit in the
proper incidence permits $E$ to be torus invariant.  The image $U$ is the
direct sum of its $\dim B$ distinct weight lines, each with coefficients
$\pm1$.  In each such weight, omit the last nonzero coordinate of its image
vector.  The remaining coordinate vectors, together with $U$, are a basis
of $A\otimes B$ over $\mathbb Z$.  This specifies all quotient coordinates
without an arithmetic exceptional prime.

For each orientation and target weight, the programs form the wedge-map
columns explicitly and eliminate the image of $U$.  If $V_w$ is a quotient
weight block and $x\in V_w$, the projected column indexed by a factor basis
vector $f_i$ is a linear map $C_{w,i}x$ of target weight $w+\operatorname{wt}(f_i)$.
These target weights are distinct for fixed $w$.  Thus adjoining a line
$kx$ increases rank by precisely the number of its nonzero projected columns.
All the $C_{w,i}$ are recomputed, rather than read from the inherited data.

The finite enumeration is as follows.
\begin{enumerate}
\item Enumerate the support of the quotient $E/U$, respecting the capacity
of every weight block.  For quotient dimension one or two, the supports
are single blocks of sufficient capacity or pairs of distinct blocks.
For quotient dimension four, enumerate every support of at most four blocks
whose capacities sum to at least four.
\item Saturate each selected block by its whole quotient weight space.
If the kernel of either wedge map on this larger space has dimension
less than $r$, the support is impossible: the kernel for $E$ injects into it.
\item For a quotient line, enumerate every possible simultaneous zero
pattern of the $C_{w,i}$.  Start with no equations, repeatedly adjoin the
row space of one $C_{w,i}$, and keep intersections with a nonzero kernel.
Canonical exact row reduction removes repetitions.  An intersection $K$
represents the stratum obtained by removing the kernels of the remaining
maps that do not vanish identically on $K$.  It is nonempty over the infinite
algebraic closure: a finite union of proper linear subspaces cannot cover
$K$.  Conversely, every vector's exact zero pattern occurs in this enumeration.
\item For lines in distinct weights, the number of distinct active target
weights is a joint rank lower bound.  Another valid bound sums the individual
increments and subtracts the total possible overlap of target weights.
For a fixed line, quotient by its actual image before testing the other line.
\item For the support-nine cases the only surviving multiplicity patterns
are $(1,1,1,1)$ and $(2,1,1)$.  In the first pattern, quotient by the images
of every fixed line.  One or two parameter lines remain.  Repeat the exact
arrangement enumeration in this quotient and retain a refined stratum only
if every formerly required nonzero map remains nonzero somewhere on it.
For a plane in a three-dimensional block, parameterize it as $\ker n$.
A column map of rank one restricts to rank zero exactly when $n$ spans its
row space; a map of rank two restricts to rank one exactly when $n$ annihilates
its one-dimensional kernel; ranks zero and three restrict to ranks zero and
two.  These are linear conditions on $n$, so the same complete arrangement
algorithm applies.  Sum, over target weights, the largest contributed rank.
\end{enumerate}
These operations exhaust all subspaces, including the parameter families.
No generic sample stands in for a projective weight space.

For the first 20 cases, \texttt{check\_twodim\_extensions.py} gives:
\[
\begin{array}{c|r|r|r|r}
&\text{supports}&\text{saturated}&\text{line tests}&\text{conditional}\\\hline
 P09A&1095&1042&53&0\\
 P09B&50&47&3&0\\
 P12C&52&50&2&0\\
 P12D&60&58&2&0\\
 P14E&1189&1185&4&0\\
 P14F&1143&1076&66&1\\
 P13.1&60&58&2&0\\
 P13.2&52&49&3&0\\
 P13.3&50&45&5&0\\
 P13.4&50&46&4&0\\
 P13.5&56&56&0&0\\
 P16.1&1189&1176&13&0\\
 P16.2&1143&1075&67&1\\
 P16.3&1190&1173&17&0
\end{array}
\]
The six determinant representatives account for twelve cases by transposition.
The two support-nine permanent types and determinant type are handled by
\texttt{check\_support9.py}; the determinant transpose is the fourth case.
Their total supports are respectively $46119,51985,46119$, with no survivors.
The exact counts, including all parameter strata, are in the output JSON files.

For the remaining P-19 case, the first two steps leave eight supports.
Exact joint refinement gives eight fixed extensions, each of dimension 13.
The higher condition is
\[
 \dim\mathcal P_{22}(E)\ge r,
 \quad
 \mathcal P_{22}(E)=
 (A\otimes E\otimes B)\cap
 ((A^{\otimes2})^{\mathfrak S_2}\otimes(B^{\otimes2})^{\mathfrak S_2}).
\]
It is again a fixed-ambient kernel condition and therefore closed.  Its
necessity follows from the $r$ independent tensors $u_i^{\otimes2}\otimes
v_i^{\otimes2}$ for a span of pure tensors, with the same proper-incidence
argument.  Using invariant tensors, including diagonal entries, is essential
in characteristic two; no division by two occurs.
The program \texttt{check\_higher\_prolongation.py} first computes the kernel
of the $385\times143$ $B$-wedge map: rank 130 and kernel dimension 13.
It embeds this kernel in $A\otimes B^{\otimes2}$, verifies symmetry by exact
coefficient equality, and forms the next $A$-wedge map.  That map has domain
dimension 91 and rank 79.  Thus $\dim\mathcal P_{22}(E)=12$ in all eight cases,
a contradiction.  These computations reconstruct the matrices from $S$;
they import no old certificates.  This completes the independent proof.

\subsection*{BR-002: forty-three in every characteristic}
For the integral divided-power tensors fixed in BR\_PENCIL FL-01 and every
algebraically closed field $k$,
\[
 \mu_7^{\det}=\mu_7^{\mathrm{per}}=14,
 \qquad
 \underline R_k(S_{\det,4})\ge43,
 \qquad
 \underline R_k(S_{\mathrm{per},4})\ge43.
\]
The assertion includes characteristics two and three.

\paragraph{Proof.}
Use the field-independent coordinate degeneration BR\_PENCIL P-01 and
its determinant compression P-07.  These are integral root-subgroup
specializations, so P-07 remains valid in characteristic two.  There are
25 determinant Borel seven-plane ideals.  For the permanent, the 1517
low-support coordinate orbits under row and column permutations and
transposition cover exactly 1543608 subsets; support at least 14 gives
a unit flattening minor.  These finite cover assertions, carried from
BR\_PENCIL R-08/R-10, are also independently rerun by
\texttt{check\_cover.py}, which exhausts all $\binom{36}{7}=8347680$ subsets.

For every remaining non-prolongation case, form the entire Koszul map,
rather than the inherited chosen minor.  If its exterior factor has dimension
$t$ and degree $p$, a pure tensor has map rank $\binom{t-1}{p}$ in every
characteristic: choose its nonzero factor as a basis vector and inspect the
basis wedges.  Rank is subadditive and its upper-bound locus is closed.
Consequently map rank $\rho$ gives border rank at least
$\lceil\rho/\binom{t-1}{p}\rceil$.
The independent checker \texttt{check\_koszul.py} forms all columns directly
from $T_S^\epsilon$, separates exact torus-character blocks, and computes
all block ranks over $\mathbb Q$.  All 11 determinant support/Koszul cases
and 1506 permanent Koszul cases give at least 14.

Here is the arithmetic justification, beyond a computation at one large prime.
Every rational elimination records the numerator and denominator of every
nonunit pivot.  Direct nonzero scalar tests in the arrangement-compatibility
step also record their numerators and denominators.  The primes that occur
divide six.  Every row operation and every inverse used by the proof is
therefore defined over $\mathbb Z[1/6]$; every exact zero identity holds there.
All rank computations, linear intersections, and their quotient maps commute
with specialization to characteristics other than two and three.  A nonzero
branch test also remains nonzero there.  Intersections that become identical
can only duplicate already enumerated strata, not create omitted ones.
This proves the full calculation uniformly away from two and three.

The exceptional computations are small and separate.
\begin{itemize}
\item In characteristic three, all eleven determinant support/Koszul cases
again give at least 14.  The two-dimensional-extension tests have no prime
three in their pivots.  The support-nine determinant arrangement is rerun
over $\mathbb F_3$ and has no survivors.  The permanent Koszul calculation
has only prime two in its pivots; its support-nine cases are also rerun over
$\mathbb F_3$ with no survivors.
\item In characteristic two, determinant and permanent tensors coincide.
The determinant compression reduces to the same 25 Borel ideals.
All eleven support/Koszul cases still give at least 14 (the two support-factor
ranks become 6457, still greater than $13\cdot495=6435$).
Five of the six ordinary determinant prolongation representatives pass the
same exhaustive computation over $\mathbb F_2$.  The P-09 A representative
has eight survivors of the first-stage test; their joint arrangements give
eight fixed extensions.  The higher checker gives $\mathcal P_{22}$ dimension
12 for each, exactly as above.  Its transpose follows by swapping rows and
columns.  The support-nine determinant tensor equals the permanent A tensor,
whose full enumeration over $\mathbb F_2$ has no survivors.
\end{itemize}
Thus every required seventh residual tensor has border rank at least 14
in every characteristic.  The integral anchored-shape curve BR\_PENCIL P-06
supplies the matching upper bound.  Finally the flattening image of the
fourfold tensor has dimension 36, with unit coefficient generators, and
BR\_PENCIL R-01 gives $36-7+14=43$.

\paragraph{What is and is not claimed about unit minors.}
This proof supplies arithmetic replacements for the rank certificates;
it does not assert that all 9038 originally selected determinants are units.
In fact a literal replacement retaining every full block rank is impossible,
as BR-003 shows.  Rank losses in small characteristics are harmless when the
remaining total rank still exceeds the required pure-tensor threshold.

\subsection*{BR-003: a genuine nonunit block in the original certificate architecture}
For the determinant subset $S=\{0,1,2,3,6,7,12\}$, take the $B$-Koszul map
of exterior degree four.  In the weight
$(2,0,0,1;1,0,0,2)$, the complete nonzero block, in the ordered bases specified
by \texttt{even\_minor\_block.json}, is
\[
 M=\begin{pmatrix}
 1&0&1&0\\1&0&0&1\\0&1&0&0\\0&0&1&1
 \end{pmatrix}.
\]
It has determinant 2 and rank 3 in characteristic two.  In particular
there is no unit four-minor in that block, nor any four-minor nonzero in
characteristic two.  The original characteristic-zero full block rank cannot
be preserved by merely changing a selected minor in this block.
The complete support-factor rank nevertheless drops only from 6501 to 6457,
which still proves the same border-rank lower bound 14.
The proof is the displayed determinant computation and the independent full
block reconstruction; its first three independent rows certify rank three
modulo two.

\subsection*{Algebraic-geometric dictionary for BR-001--BR-003}
All factors are finite free modules over $\mathbb Z$.  Wedge maps are morphisms
between their exterior-power bundles, and their rank conditions define
closed determinantal loci.  Candidate extension spaces form closed incidences
in relative Grassmannians.  Projectivity supplies proper projection and
cocharacter specialization.  The arithmetic result proves emptiness in each
algebraically closed fiber: first over $\operatorname{Spec}\mathbb Z[1/6]$
by invertible elimination, then separately over the fibers at two and three.
It does not infer emptiness of special fibers from emptiness of the generic
fiber.  Invariant symmetric tensor squares in the higher test are free
integral modules, so the proof also specifies the characteristic-two model.

The inherited exact foundational locators are EGA II, Definition 5.4.1,
Proposition 5.4.2, and Theorem 5.5.3(i) (properness and projective morphisms);
FGA, Bourbaki expos\'e 221, Theorems 3.1--3.2 (Quot/Hilbert parameter spaces);
SGA 3, Expos\'e I, Sections 4.5--4.6 including Definition 4.6.1 (module and
linear group functors); and SGA 6, Expos\'e V, Definition 2.1, Examples 2.2(b),
and Lemma 2.2.1 (exterior operations on $K_0$).  These are framework locators
carried from BR\_PENCIL's checked dictionary, not fresh facsimile inspections
and not attributions of the modern numerical bounds to Grothendieck.
No unpublished archive cote/scan for these findings has been established.
No purity, weight, $\ell$-adic, or motivic statement is used.
